\documentclass[11pt]{article}
\usepackage[margin=1in]{geometry}
\usepackage{amsmath,amssymb,mathtools}
\usepackage{enumitem}
\usepackage{booktabs}
\usepackage{microtype}
\usepackage[hidelinks]{hyperref}
\usepackage{xurl}
\setlength{\parindent}{0pt}
\setlength{\parskip}{0.55em}
\newcommand{\E}{\mathbb{E}}
\newcommand{\1}{\mathbf{1}}

\title{A Two-Period Directed-Search Model with Wealth Heterogeneity\\
\large Equilibrium Definition and Numerical Solution}
\author{}
\date{September 2026}

\begin{document}
\maketitle

\section{Environment}

This note studies a simplified version of the two-period model in Repele (2026). The simplification removes heterogeneity in job type: all jobs have the same productivity, separation probability, and vacancy-posting cost. Workers differ only in initial liquid wealth. The remaining source of heterogeneity across submarkets is the posted wage, which is associated in equilibrium with a different market tightness and hence a different job-finding probability.

There are two periods, $t=1,2$. A continuum of workers is indexed by initial liquid wealth
\[
a_1 \sim G(a_1), \qquad a_1 \in [\underline a_1,\overline a_1].
\]
Workers have expected utility
\[
u(c_1)+\beta\E[u(c_2)],
\]
where $u'(c)>0$ and $u''(c)<0$. They can save in a risk-free asset with gross return $R$ and face the borrowing constraint
\[
a_2\geq \underline a.
\]
All workers start period 1 unemployed. If a worker is unemployed in period 2, she receives income $b$.

All jobs have common productivity $\bar y$, common separation probability $\delta$, and common vacancy-posting cost $\phi$. Define the retention probability
\[
\Delta \equiv 1-\delta.
\]

A submarket is indexed by its posted wage $w$. Let $n(w)$ denote the mass of workers applying to submarket $w$ and $v(w)$ the mass of vacancies. Market tightness is
\[
\theta(w)=\frac{v(w)}{n(w)}
\]
for an active submarket. A worker finds a job with probability $p(\theta)$, where $p'(\theta)>0$, while a vacancy is filled with probability
\[
q(\theta)=\frac{p(\theta)}{\theta}, \qquad q'(\theta)<0.
\]

The period-2 values are
\[
E_2(a_2,w)=u(Ra_2+w), \qquad U_2(a_2)=u(Ra_2+b),
\]
and the value of a surviving filled job is
\[
J(w)=\bar y-w.
\]

\section{Equilibrium}

A directed-search equilibrium is a collection
\[
\mathcal E=\left\{a_2^*(a_1),\,w^*(a_1),\,\theta(w),\,n(w),\,v(w)\right\}
\]
such that the following conditions hold.

\begin{enumerate}[label=\textbf{\arabic*.},leftmargin=*]

\item \textbf{Worker optimization.} Taking the tightness schedule $\theta(w)$ as given, every worker with initial wealth $a_1$ solves
\[
\begin{split}
\left(a_2^*(a_1),w^*(a_1)\right) \in
\arg\max_{a_2\geq \underline a,\,w}\; &u(Ra_1-a_2) \\
&+\beta p(\theta(w))\Delta\,u(Ra_2+w)\\
&+\beta\left[1-p(\theta(w))\Delta\right]u(Ra_2+b).
\end{split}
\tag{W}
\]
The probability that the worker is employed in period 2 is $p(\theta(w))\Delta$: the worker must first match and the newly created match must then survive.

\item \textbf{Firm free entry.} For every potential wage submarket $w$, the value of a vacancy is
\[
V(w)=-\phi+\beta\Delta q(\theta(w))(\bar y-w).
\]
Free entry and complementary slackness require
\[
\phi \geq \beta\Delta q(\theta(w))(\bar y-w), \qquad \theta(w)\geq 0,
\]
\[
\theta(w)\left[\phi-\beta\Delta q(\theta(w))(\bar y-w)\right]=0.
\tag{FE}
\]
Thus, whenever $\theta(w)>0$,
\[
\phi=\beta\Delta q(\theta(w))(\bar y-w).
\]

\item \textbf{Worker-flow consistency.} Workers' optimal search choices determine the mass of applicants to each submarket:
\[
n(w)=\int \1\{w^*(a_1)=w\}\,dG(a_1).
\tag{C1}
\]
For a discrete wealth approximation with masses $\mu_i$, this becomes
\[
n_j=\sum_i \mu_i\,\1\{w_i^*=w_j\}.
\]

\item \textbf{Vacancy/tightness consistency.} For every active submarket,
\[
v(w)=\theta(w)n(w).
\tag{C2}
\]
Hence $\theta(w)=v(w)/n(w)$ whenever $n(w)>0$. If no worker chooses a submarket, set $n(w)=v(w)=0$; the zero-profit schedule $\theta(w)$ can still be viewed as the tightness that would prevail were that submarket active.

\item \textbf{Matching consistency.} Job-finding and vacancy-filling probabilities are generated by the same tightness:
\[
p_w=p(\theta(w)), \qquad q_w=q(\theta(w))=\frac{p(\theta(w))}{\theta(w)}.
\tag{C3}
\]

\end{enumerate}

There is no asset-market-clearing condition because $R$ is taken as exogenous.

\section{Why there is no fixed point over the tightness schedule}

In this simplified economy, conditional on a wage $w$, free entry determines $\theta(w)$ without reference to the mass or composition of workers who choose that submarket. The firm's expected profit depends on $w$ and $\theta$, but not on the absolute numbers of workers and vacancies. Constant returns to scale in matching imply that only their ratio matters.

For an active submarket,
\[
q(\theta(w))=\frac{\phi}{\beta\Delta(\bar y-w)}.
\tag{1}
\]
Since $q(\theta)$ is decreasing, this determines a unique $\theta(w)$ whenever the right-hand side lies in the range of $q$. Workers then choose among the resulting submarkets. Their choices determine $n(w)$, and vacancy entry adjusts mechanically to
\[
v(w)=\theta(w)n(w).
\]
Thus the computational sequence is
\[
w \longrightarrow \theta(w) \longrightarrow \text{worker choices} \longrightarrow n(w) \longrightarrow v(w).
\]
There is no need to guess $\theta(w)$, solve workers' choices, and iterate back to a new $\theta(w)$. A genuine fixed point would arise if firm profitability depended on the endogenous composition of applicants, for example if worker productivity differed by wealth or another worker type.

\section{A convenient matching-function parameterization}

Following the functional form used by Repele (2026, Appendix B1), consider the constant-returns matching function
\[
M(n,v)=nv\left(n^\gamma+v^\gamma\right)^{-1/\gamma}, \qquad \gamma>0.
\]
With $\theta=v/n$, the worker job-finding and firm vacancy-filling probabilities are
\[
p(\theta)=\theta\left(1+\theta^\gamma\right)^{-1/\gamma},
\]
\[
q(\theta)=\left(1+\theta^\gamma\right)^{-1/\gamma}.
\]
These probabilities are bounded between zero and one and satisfy $q(\theta)=p(\theta)/\theta$.

For this matching function, the free-entry schedule can be computed without a root finder. Define
\[
q^{FE}(w)\equiv \frac{\phi}{\beta\Delta(\bar y-w)}.
\]
For $0<q^{FE}(w)<1$,
\[
\theta(w)=\left(\left[q^{FE}(w)\right]^{-\gamma}-1\right)^{1/\gamma},
\tag{2}
\]
and then
\[
p(w)=p(\theta(w)).
\]
If $q^{FE}(w)\geq 1$, the vacancy cannot cover its posting cost even when it is filled with probability one, so that submarket is inactive and $\theta(w)=0$.

Repele's quantitative model uses $\gamma=0.38$, which implies an average elasticity of job finding with respect to tightness in the range targeted in her calibration. For the simplified two-period exercise, $\gamma=0.38$ is therefore a natural benchmark, although it is not necessary for the theoretical results.

\section{Numerical parameterization and grids}

A useful way to parameterize the toy model is to normalize productivity and calibrate the vacancy cost to a reference job-finding rate rather than choosing $\phi$ arbitrarily.

\subsection*{Preferences and common job parameters}

A convenient utility specification is CRRA:
\[
u(c)=\frac{c^{1-\sigma}}{1-\sigma}, \qquad \sigma>0,\;\sigma\neq 1,
\]
with $u(c)=\log c$ when $\sigma=1$. Repele's quantitative model uses $\sigma=2$, which is a useful benchmark here as well.

Normalize
\[
\bar y=1.
\]
Choose $R$ and $\beta$ consistently with the interpretation of one model period. Choose the common separation probability $\delta$ from the desired time frequency, and set $\Delta=1-\delta$. A simple benefit parameterization is
\[
b=\rho_b\bar y,
\]
where $\rho_b$ is a replacement-rate parameter.

\subsection*{Calibrating the vacancy cost}

Choose a reference wage $\bar w$ and a target job-finding probability $\bar p\in(0,1)$. Under the matching function above,
\[
\bar\theta=\left(\frac{\bar p^\gamma}{1-\bar p^\gamma}\right)^{1/\gamma},
\]
and
\[
\bar q=q(\bar\theta)=\left(1-\bar p^\gamma\right)^{1/\gamma}.
\]
Then set
\[
\phi=\beta\Delta\bar q(\bar y-\bar w).
\tag{3}
\]
This guarantees that $(\bar w,\bar p)$ lies on the free-entry schedule.

\subsection*{Wage grid}

For an active submarket the free-entry equation requires
\[
q^{FE}(w)\leq 1,
\]
which implies
\[
w\leq \bar y-\frac{\phi}{\beta\Delta}.
\]
Moreover, with no nonpecuniary value of employment, wages weakly below $b$ are dominated for workers. A practical wage grid is therefore
\[
w\in[w_{\min},w_{\max}],
\]
with
\[
w_{\min}=b+\varepsilon_w,
\]
\[
w_{\max}=\bar y-\frac{\phi}{\beta\Delta}-\varepsilon_w,
\]
where $\varepsilon_w>0$ is small. The parameterization should satisfy $w_{\min}<w_{\max}$. A grid of roughly $N_w=200$--$500$ points is more than sufficient for this two-period problem.

A uniform wage grid is transparent. Numerically, however, $\theta(w)$ can change quickly near the upper wage bound. An alternative is to use a grid that is uniform in $p$ or $\theta$ and recover the corresponding wage from
\[
w(\theta)=\bar y-\frac{\phi}{\beta\Delta q(\theta)}.
\]

\subsection*{Wealth grid}

With a borrowing limit $\underline a=0$, choose
\[
a_1\in[0,a_{\max}].
\]
It is useful to place more grid points near zero because wealth has the largest effect on search choices close to the borrowing constraint. For example, with $i=1,\ldots,N_a$,
\[
a_{1,i}=a_{\max}\left(\frac{i-1}{N_a-1}\right)^\kappa,
\qquad \kappa>1.
\]
Values such as $\kappa=2$ and $N_a=200$ are adequate for a simple numerical illustration. Attach probability weights $\mu_i$ either from a chosen parametric distribution or directly from empirical liquid-wealth data.

\section{Solution algorithm}

For a discrete numerical implementation, the following algorithm computes the equilibrium directly.

\begin{enumerate}[label=\textbf{Step \arabic*.},leftmargin=*]

\item Choose grids $\{w_j\}_{j=1}^{N_w}$ and $\{a_{1,i}\}_{i=1}^{N_a}$ and probability weights $\{\mu_i\}$ for the wealth distribution.

\item For every wage $w_j$, compute
\[
q_j^{FE}=\frac{\phi}{\beta\Delta(\bar y-w_j)}.
\]
If $0<q_j^{FE}<1$, compute $\theta_j$ from equation (2), and then $p_j=p(\theta_j)$. If $q_j^{FE}\geq 1$, set the submarket to inactive.

\item For each pair $(a_{1,i},w_j)$, solve the worker's one-dimensional savings problem
\[
\max_{a_2\geq \underline a}
\left\{
u(Ra_{1,i}-a_2)
+\beta p_j\Delta u(Ra_2+w_j)
+\beta(1-p_j\Delta)u(Ra_2+b)
\right\}.
\tag{4}
\]
With positive-consumption utility and $\underline a\geq0$, a convenient feasible interval is
\[
a_2\in[\underline a,\,Ra_{1,i}-\varepsilon_c].
\]
The problem can be solved by a scalar optimizer or, when the solution is interior, by the Euler equation.

\item For each wealth type $a_{1,i}$, compare the optimized value across wage submarkets and choose
\[
j^*(i)\in\arg\max_j \widetilde U(a_{1,i};w_j).
\]
Set
\[
w_i^*=w_{j^*(i)}, \qquad \theta_i^*=\theta_{j^*(i)},
\]
and retain the corresponding optimal savings $a_{2,i}^*$.

\item Count the mass of workers applying to each submarket:
\[
n_j=\sum_i \mu_i\1\{j^*(i)=j\}.
\]
This is the mass of period-1 searchers in submarket $j$, not an unemployment rate.

\item For every active submarket, create the mass of vacancies
\[
v_j=\theta_j n_j.
\]
If $n_j=0$, set $v_j=0$.

\item Verify equilibrium numerically. For every active submarket check
\[
\left|\phi-\beta\Delta q(\theta_j)(\bar y-w_j)\right|<\text{tol},
\]
and
\[
\left|\theta_j-\frac{v_j}{n_j}\right|<\text{tol}.
\]
Also verify that each worker's chosen wage attains the maximum of the worker problem over the wage grid.

\end{enumerate}

The output of the algorithm is the wealth-dependent policy map
\[
a_1\mapsto \left(a_2^*(a_1),w^*(a_1),\theta^*(a_1),p^*(a_1)\right),
\]
together with the equilibrium masses $n(w)$ and $v(w)$. The main economic object of interest is how wealth changes the worker's willingness to trade a higher wage for a lower job-finding probability.

\section{Reference}

\begin{thebibliography}{9}
\bibitem{Repele2026}
Repele, Amalia (2026), \emph{Wealth Sorting and Cyclical Employment Risk}, Job Market Paper, International Institute for Economic Studies, Stockholm University, July 2026.
\end{thebibliography}

\end{document}